\clearpage

\chapter{SUPPORTING LITERATURE}

\chapter{Literature Review}
A systematic review of contemporary literature was conducted to evaluate existing methodological frameworks, feature engineering representations, classification algorithms, class imbalance mitigation techniques, and dataset challenges in clinical natural language processing. The review focuses on peer-reviewed research analyzing medical narratives, electronic health records, and document categorization systems.

\chapter{Paper 1 -- Omar AM, Elhoseny M, Hassanien AM. Clinical text classification with word representation features and machine learning algorithms. International Journal of Online and Biomedical Engineering (iJOE). 2023.}

The paper titled ``Clinical Text Classification with Word Representation Features and Machine Learning Algorithms'' \cite{omar2023} investigates the impact of diverse text representation techniques on Electronic Medical Record (EMR) narratives. The authors introduced a structured four-phase classification framework consisting of clinical text preprocessing, word vectorization, dimensionality reduction, and supervised classification.

To model clinical narratives, the study conducted a comparative assessment of classical Bag-of-Words (BoW) and Term Frequency--Inverse Document Frequency (TF-IDF) representations against dense distributed semantic embeddings (Word2Vec) across multiple supervised estimators, including Logistic Regression, Support Vector Machine (SVM), Naive Bayes, and k-Nearest Neighbors (k-NN). The empirical evaluations indicated that pairing continuous Word2Vec embeddings with a k-NN classifier yielded the highest classification accuracy of 92\%, demonstrating that continuous representations capture semantic affinities between medical terms more effectively than unigram Bag-of-Words.

However, the authors noted that distance-based k-NN introduces significant computational latency during test-time inference, scaling poorly to large-scale hospital databases. Furthermore, the investigation evaluated individual classifiers in isolation without exploring multi-model ensemble consensus or real-time web deployment frameworks. Crucially, Principal Component Analysis (PCA) was evaluated exclusively as an offline dimensionality reduction step in Paper 1 and is not required in our proposed architecture, as sublinear TF-IDF with feature pruning maintains optimal sparsity without dense projection.

\begin{table}[H]
\centering
\small
\setlength{\tabcolsep}{5pt}
\begin{tabularx}{\linewidth}{|>{\raggedright\arraybackslash}p{3.8cm}|>{\raggedright\arraybackslash}X|}
\hline
\textbf{TITLE} & Clinical Text Classification with Word Representation Features and Machine Learning Algorithms \cite{omar2023} \\
\hline
\textbf{AREA OF WORK} & Clinical Natural Language Processing / EMR Text Vectorization \\
\hline
\textbf{DATASET} & Electronic Medical Record (EMR) Clinical Transcriptions Dataset \\
\hline
\textbf{METHODOLOGY / STRATEGY} & 4-Phase Framework: (1) Preprocessing, (2) Vectorization (BoW, TF-IDF, Word2Vec), (3) Feature Reduction (PCA), (4) Supervised Classification \\
\hline
\textbf{ALGORITHM} & Logistic Regression, Support Vector Machine (SVM), Naive Bayes, k-NN, PCA \\
\hline
\textbf{RESULT / PRECISION} & Word2Vec + k-NN achieved the highest classification accuracy of 92\% \\
\hline
\textbf{ADVANTAGES} & Proves dense semantic representations capture clinical context better than raw unigram BoW \\
\hline
\textbf{LIMITATIONS} & Distance-based k-NN exhibits high inference latency; lacks ensemble voting, open-set handling, and real-time web deployment \\
\hline
\textbf{RELEVANCE TO MY PROJECT} & Validates the standard 4-phase clinical text preprocessing sequence and benchmark machine learning estimators \\
\hline
\textbf{FUTURE PROPOSAL} & Develop lightweight ensemble classifiers suitable for instantaneous CPU web inference with calibrated confidence \\
\hline
\end{tabularx}
\caption{Summary of Paper 1}
\label{tab:summary_paper1}
\end{table}

\chapter{Paper 2 -- Zhang H, Zhu D, Tan H, Shafiq M, Gu Z. Medical specialty classification based on semiadversarial data augmentation. Hindawi Computational Intelligence and Neuroscience. 2023.}

The research article titled ``Medical Specialty Classification Based on Semiadversarial Data Augmentation'' \cite{zhang2023} investigated extreme category imbalance and data sparsity in clinical specialty categorization. The authors introduced Semiadversarial Data Augmentation (SemiADA), a framework generating synthetic clinical text samples along adversarial perturbation trajectories to populate sparse decision regions in high-dimensional feature spaces.

Additionally, a Probabilistic Information (PI) layer was incorporated following the Softmax output to recalculate prediction confidence by weighting domain-specific clinical nouns. The authors evaluated a filtered subset of the Kaggle Medical Specialty Classification dataset restricted to 18 classes, achieving an average performance improvement of 14.9\% in accuracy and F1-score over baseline benchmarks. Their findings confirmed that domain-specific clinical nouns carry the primary discriminative signal for medical department routing.

Despite these gains, generating synthetic samples along adversarial attack paths incurs extreme computational overhead. Furthermore, fine-tuning deep transformer architectures (such as BERT or BioBERT) requires dedicated GPU clusters and introduces substantial inference latency, creating a major barrier for real-time clinical web applications operating in resource-constrained hospital environments. The study also operated strictly under a closed-set assumption, omitting mechanisms to intercept out-of-domain records.

\begin{table}[H]
\centering
\small
\setlength{\tabcolsep}{5pt}
\begin{tabularx}{\linewidth}{|>{\raggedright\arraybackslash}p{3.8cm}|>{\raggedright\arraybackslash}X|}
\hline
\textbf{TITLE} & Medical Specialty Classification Based on Semiadversarial Data Augmentation \cite{zhang2023} \\
\hline
\textbf{AREA OF WORK} & Clinical Text Data Augmentation / Severe Class Imbalance Mitigation \\
\hline
\textbf{DATASET} & Kaggle Medical Specialty Classification Dataset (Subset restricted to 18 classes) \\
\hline
\textbf{METHODOLOGY / STRATEGY} & Semiadversarial sample generation along attack trajectories + Probabilistic Information (PI) noun-weighted Softmax recalculation \\
\hline
\textbf{ALGORITHM} & SemiADA, BERT (Transformer), Softmax with PI Noun Layer \\
\hline
\textbf{RESULT / PRECISION} & Achieved an average 14.9\% boost in Accuracy and Macro F1-score over baseline benchmarks \\
\hline
\textbf{ADVANTAGES} & Populates sparse minority decision boundaries using geometry-aware synthetic clinical samples \\
\hline
\textbf{LIMITATIONS} & Extreme training latency; fine-tuning BERT demands dedicated GPU accelerators; closed-set classification only \\
\hline
\textbf{RELEVANCE TO MY PROJECT} & Highlights class imbalance and noun significance in MTSamples data; motivates high-dimensional TF-IDF with ensemble ML \\
\hline
\textbf{FUTURE PROPOSAL} & Combine sublinear TF-IDF n-gram weighting with probability-calibrated ensemble learning and open-set rejection on standard CPUs \\
\hline
\end{tabularx}
\caption{Summary of Paper 2}
\label{tab:summary_paper2}
\end{table}

\chapter{Paper 3 -- Blanchard AE, Gao S, Yoon HJ, Tourassi G. A keyword-enhanced approach to handle class imbalance in clinical text classification. IEEE Journal of Biomedical and Health Informatics. 2022.}

Blanchard et al., writing in the \textit{IEEE Journal of Biomedical and Health Informatics} \cite{blanchard2022}, addressed the challenge of deep neural networks overfitting on rare clinical categories when trained on lengthy, complex pathology narratives. The authors introduced a data-centric regularisation strategy by injecting short class-associated keyword sequences into mini-batches during neural network training, formulating a dual loss function:
\begin{equation}
L = L_{\text{docs}} + \alpha L_{\text{key}}
\end{equation}
where $L_{\text{docs}}$ evaluates cross-entropy loss over full documents, $L_{\text{key}}$ penalizes misclassifications on isolated diagnostic keyword sequences, and $\alpha$ balances the objective.

Evaluating on the National Cancer Institute (NCI) Surveillance, Epidemiology, and End Results (SEER) cancer pathology registry dataset, the study benchmarked Convolutional Neural Networks (CNNs), keyword-constrained loss functions, and conventional resampling techniques (oversampling and undersampling). The findings demonstrated that keyword constraints dramatically improved Macro F1-score and minority class recall without degrading majority class performance. Crucially, the authors emphasized that overall micro-accuracy is an inherently deceptive evaluation metric in clinical NLP, as high global accuracy frequently masks near-zero recall on rare clinical conditions.

However, their approach relies on manual curation of comprehensive keyword lexicons by clinical domain experts for every target category, limiting rapid scalability to new medical specialties. This limitation directly motivates our system's automated statistical TF-IDF approach, which extracts informative unigram and bigram vocabulary automatically without requiring hand-crafted dictionaries.

\begin{table}[H]
\centering
\small
\setlength{\tabcolsep}{5pt}
\begin{tabularx}{\linewidth}{|>{\raggedright\arraybackslash}p{3.8cm}|>{\raggedright\arraybackslash}X|}
\hline
\textbf{TITLE} & A Keyword-Enhanced Approach to Handle Class Imbalance in Clinical Text Classification \cite{blanchard2022} \\
\hline
\textbf{AREA OF WORK} & Imbalanced Clinical Narrative Categorization / Pathology NLP \\
\hline
\textbf{DATASET} & NCI SEER Surveillance Cancer Pathology Registry Reports Dataset \\
\hline
\textbf{METHODOLOGY / STRATEGY} & Mini-batch keyword sequence injection to constrain the neural loss function and prevent rare-class forgetting \\
\hline
\textbf{ALGORITHM} & Convolutional Neural Networks (CNNs), Keyword-Constrained Loss Optimization \\
\hline
\textbf{RESULT / PRECISION} & Significantly boosted Macro F1-score and minority recall across imbalanced pathology classes \\
\hline
\textbf{ADVANTAGES} & Demonstrates that keyword and n-gram terms carry decisive clinical signals; confirms Macro F1 as the primary clinical metric \\
\hline
\textbf{LIMITATIONS} & Requires exhaustive manual curation of keyword dictionaries by medical specialists; computationally heavy neural training \\
\hline
\textbf{RELEVANCE TO MY PROJECT} & Proves that n-gram keyword weighting preserves rare class discrimination; validates Macro F1 as our primary evaluation metric \\
\hline
\textbf{FUTURE PROPOSAL} & Implement automated statistical TF-IDF unigram and bigram extraction with sublinear scaling, eliminating manual dictionary construction \\
\hline
\end{tabularx}
\caption{Summary of Paper 3}
\label{tab:summary_paper3}
\end{table}

\chapter{Literature Review Summary}
Table \ref{tab:overall_summary} presents a comparative synthesis of the reviewed literature alongside our proposed MedSpecialty AI Platform.

\begin{table}[H]
\centering
\small
\setlength{\tabcolsep}{4pt}
\begin{tabularx}{\linewidth}{|>{\raggedright\arraybackslash}p{2.2cm}|>{\raggedright\arraybackslash}p{2.1cm}|>{\raggedright\arraybackslash}p{2.1cm}|>{\raggedright\arraybackslash}X|>{\raggedright\arraybackslash}p{2.2cm}|>{\raggedright\arraybackslash}X|}
\hline
\textbf{Paper} & \textbf{Area} & \textbf{Dataset} & \textbf{Algorithm} & \textbf{Advantages} & \textbf{Relation to Proposed Work} \\
\hline
Omar et al. \cite{omar2023} & Clinical Text Vectorization & EMR Text & LR, SVM, NB, k-NN & High accuracy (92\%) with Word2Vec & Establishes the 4-phase text preprocessing pipeline \\
\hline
Zhang et al. \cite{zhang2023} & Data Augmentation & MTSamples (18 classes) & SemiADA, BERT & +14.9\% gain in Accuracy and F1 & Identifies MTSamples class imbalance and clinical noun signal \\
\hline
Blanchard et al. \cite{blanchard2022} & Class Imbalance Regularization & NCI SEER Pathology & CNN + Keyword Loss & High Macro F1 on minority classes & Establishes Macro F1 as primary metric; validates n-gram term weighting \\
\hline
\textbf{Proposed Work} & \textbf{Specialty Triage \& Open-Set Routing} & \textbf{MTSamples (Curated 8 Core + Other)} & \textbf{Sublinear TF-IDF (8k) + Soft Voting Ensemble [2:2:1:1]} & \textbf{Fast CPU inference, open-set safety gate ($\tau=0.48$), interactive web UI} & \textbf{Complete, verified clinical decision support platform (MedSpecialty AI)} \\
\hline
\end{tabularx}
\caption{Overall Summary of Literature and Relation to Proposed System}
\label{tab:overall_summary}
\end{table}

\chapter{Evolution of Research}
Clinical natural language processing has evolved from elementary Bag-of-Words term frequency counters toward distributed word embeddings (Word2Vec), deep convolutional architectures, complex transformer models (BERT, BioBERT), and adversarial data augmentation (SemiADA). While deep neural networks deliver high semantic abstraction, their steep computational overhead, large parameter footprints, and GPU hardware prerequisites prevent real-time deployment in cost-conscious clinical web environments. Furthermore, deep neural networks often suffer from overconfidence on out-of-domain text.

\chapter{Identified Research Gaps}
The literature review reveals four fundamental research gaps:
\begin{enumerate}
    \item \textit{Excessive Computational Overhead}: Existing deep learning systems rely on GPU hardware accelerators and introduce substantial inference latency, making real-time interactive dictation triage challenging.
    \item \textit{Closed-Set Vulnerability}: Prior works operate exclusively under a closed-set assumption, failing to handle out-of-scope medical specialties or non-clinical narratives, resulting in hazardous false specialty assignments.
    \item \textit{Lack of End-to-End Clinical Deployment}: Most studies focus solely on offline experimental benchmarks without delivering an operational, accessible web application supporting direct PDF ingestion.
    \item \textit{Single Estimator Bias}: Many prior systems rely on a single classifier architecture, failing to exploit ensemble diversity to balance high-margin boundaries with calibrated probabilistic spreads.
\end{enumerate}

\chapter{Motivation for Proposed System}
Our project directly addresses these research gaps by designing the \textbf{MedSpecialty AI Platform}. The system couples an 8,000-dimensional sublinear TF-IDF vectorizer (capturing unigram terms and bigram clinical phrases) with a calibrated Soft Voting Ensemble aggregating Linear SVM, Logistic Regression, Random Forest, and Multinomial Naive Bayes using optimal weights [2:2:1:1]. Crucially, an Open-Set Confidence Rejection Gate ($\tau = 0.48$) is established to route unsupported or ambiguous narratives to a dedicated \textit{``Other / Requires Review''} class. The entire architecture is deployed through a responsive Flask web application featuring PyMuPDF in-memory text parsing and real-time NLP pipeline inspection on standard CPU hardware.

\chapter{Findings and Proposals}

\chapter{Findings}
The systematic review of contemporary literature yields several critical findings:
\begin{enumerate}
    \item \textit{Efficacy of Classical ML with N-Gram Term Weighting}: Classical machine learning algorithms (Support Vector Machines, Logistic Regression, Random Forest, and Naive Bayes) demonstrate competitive classification performance on clinical narratives when paired with sublinear TF-IDF feature representations, matching or exceeding heavy neural models on moderate-sized corpora while executing in milliseconds on standard CPUs.
    \item \textit{Sublinear TF Dampening Prevents Frequency Bias}: Clinical dictations feature frequent repetitions of routine procedural terms. Sublinear logarithmic scaling ($1 + \log(\text{TF})$) successfully prevents high-frequency terms from overshadowing specialty-defining diagnostic keywords.
    \item \textit{Complementary Estimator Inductive Biases}: Different classifier families exhibit complementary strengths. Linear SVM identifies robust maximum-margin hyperplanes in high-dimensional sparse spaces; Logistic Regression provides well-calibrated odds; Random Forest captures non-linear feature interactions; and Multinomial Naive Bayes contributes probabilistic frequency baselines.
    \item \textit{Necessity of Open-Set Rejection}: Closed-set classification in clinical environments creates unacceptable diagnostic risks when processing out-of-scope conditions. An empirical confidence threshold provides an effective safeguard for triaging unrecognized narratives to human experts.
    \item \textit{Importance of Macro F1 Evaluation}: In clinical NLP with imbalanced specialty distributions, overall accuracy is deceptive. Macro F1-score---which weights each clinical department equally---serves as the primary indicator of balanced diagnostic capability.
\end{enumerate}

\chapter{Proposal}
Synthesizing these findings, this project proposes the \textbf{MedSpecialty AI Platform} for automated clinical specialty classification and open-set triage:
\begin{enumerate}
    \item \textit{Curated Benchmark Dataset}: Filter the raw MTSamples corpus down from 40 overlapping/administrative categories to a curated benchmark of 1,663 records across eight mutually exclusive hospital specialties: Cardiovascular / Pulmonary, Orthopedic, Gastroenterology, Neurology, Urology, Obstetrics / Gynecology, ENT -- Otolaryngology, and Ophthalmology.
    \item \textit{Modular NLP Preprocessing}: Implement an automated cleaning pipeline performing lowercasing, regex noise removal, tokenization, stopword elimination, and WordNet morphological lemmatization.
    \item \textit{Sublinear TF-IDF Vectorization}: Extract unigram and bigram features bounded at 8,000 dimensions with sublinear term frequency scaling, document frequency cutoffs ($\text{min\_df}=2, \text{max\_df}=0.70$), and Euclidean $L_2$ normalization.
    \item \textit{Calibrated Soft Voting Ensemble}: Train Linear SVM (calibrated via Platt scaling), Logistic Regression, Random Forest, and Multinomial Naive Bayes, synthesizing their predictions using calibrated soft voting weights ($2 \times \text{SVM} + 2 \times \text{LR} + 1 \times \text{RF} + 1 \times \text{MNB}) / 6$.
    \item \textit{Open-Set Rejection Architecture}: Formulate a decision rule: if the peak ensemble probability $P_{\max} < 0.48$, the document is classified as \textit{``Other''} with status \textit{``Out of Scope / Requires Review''}, displaying nearest candidate specialties for manual triage.
    \item \textit{Interactive Web Application}: Deploy a production-grade web application featuring in-memory PDF extraction (PyMuPDF), editable clinical dictation areas, interactive TF-IDF feature tables, and multi-model consensus dashboards.
\end{enumerate}
